% Generated by tables/build_supplementary_table_s5.py from the local LOC-01 CSV.
\begin{table}[H]
  \caption{PET-ring-normalized local-geometry fractions for TfCut1 F210 and HiC L138.}
  \label{tab:per_ring_contact_geometry}
  \centering
  \small
  \setlength{\tabcolsep}{4pt}
  \renewcommand{\tabularxcolumn}[1]{m{#1}}
  \begin{tabularx}{\linewidth}{@{}>{\raggedright\arraybackslash}m{1.5cm}>{\raggedright\arraybackslash}Xccccc@{}}
    \toprule
    Target & Per-ring metric & $n_{10}/n_{20}$ & \shortstack{Mean\\L10} &
    \shortstack{Mean\\L20} & $\Delta$ & \shortstack{95\% interval\\for $\Delta$} \\
    \midrule
    TfCut1 F210 & Aromatic-ring proximity & 6/6 & 0.0343 & 0.0344 & $+0.0001$ & $[-0.0363, +0.0322]$ \\
    TfCut1 F210 & Pi-like geometry & 6/6 & 0.0067 & 0.0078 & $+0.0011$ & $[-0.0085, +0.0090]$ \\
    HiC L138 & Nonpolar contact, 0.40 nm & 5/8 & 0.0063 & 0.0105 & $+0.0041$ & $[-0.0089, +0.0172]$ \\
    HiC L138 & Nonpolar contact, 0.45 nm$^{\dagger}$ & 5/8 & 0.0097 & 0.0175 & $+0.0079$ & $[-0.0123, +0.0288]$ \\
    HiC L138 & Nonpolar contact, 0.50 nm$^{\dagger}$ & 5/8 & 0.0118 & 0.0221 & $+0.0103$ & $[-0.0143, +0.0356]$ \\
    \bottomrule
  \end{tabularx}

  \vspace{4pt}
  \parbox{\linewidth}{\footnotesize
  For each analyzed frame, the number of PET aromatic rings meeting the stated
  criterion was divided by the total number of PET aromatic rings (10 for L10
  and 20 for L20). These framewise fractions were averaged over 20--100 ns
  within each retained trajectory, and the trajectory means were averaged
  within each enzyme--length group, including zero values. $n_{10}$ and
  $n_{20}$ denote retained-trajectory counts, and $\Delta$ is the L20-minus-L10
  difference in group means. The 95\% intervals are percentile intervals from
  20,000 bootstrap resamples of whole trajectories performed separately within
  the L10 and L20 groups. Aromatic-ring proximity requires a protein--PET
  ring-centroid distance below 0.55 nm. Pi-like geometry additionally requires
  a projected ring-center offset below 0.20 nm and ring planes within
  $30\si{\degree}$ of parallel or perpendicular. The projected offset is the
  smaller projection of the centroid-to-centroid vector onto the two ring
  planes. Nonpolar contacts use the minimum distance between L138 side-chain
  carbon atoms and PET aromatic-ring carbon atoms. The 0.40-nm cutoff is
  primary, and daggers mark the 0.45- and 0.50-nm sensitivity cutoffs.
  The candidates were selected after the residue-contact screen, and the
  intervals are descriptive and unadjusted for multiple comparisons.
  Values are rounded to four decimal places from the full-precision source data.}
\end{table}
